%!%
% Copyright (c) 2026 mingcheng <mingcheng@outlook.com>
% 
% File: ch09-prompt-examples.tex
% Author: mingcheng <mingcheng@outlook.com>
% File Created: 2026-03-27 15:50:02
% 
% Modified By: mingcheng <mingcheng@outlook.com>
% Last Modified: 2026-03-29 07:50:36
%%

% !TeX root = ../prompt-engineering-guide-for-beginners.tex

\chapter{提示词示例：拿来即用的场景模板}
\label{ch:prompt-examples}

前面几章讲了很多方法论，这一章直接上「干货」。
针对日常工作和生活中最常见的场景，我准备了一批可以直接复制使用的提示词模板。
你只需要把里面的具体内容替换成你自己的需求，就能拿来即用。

\section{文本概括与总结}

日常工作中最常见的场景之一就是「读了一大堆东西，但时间不够用」。无论是会议纪要、行业报告还是长篇文章，AI 都能帮你快速抓住重点。以下几个模板覆盖了最常见的总结需求。

\subsection*{长文章一句话摘要}

\begin{promptbox}
请用一句话概括以下文章的核心观点，要求简洁精炼，不超过 50 个字：

\verb|"""|

（在这里粘贴你的文章内容）

\verb|"""|
\end{promptbox}

\subsection*{会议记录提炼行动项}

\begin{promptbox}
以下是今天的会议记录。请帮我提炼出所有的待办事项（Action Items），以编号列表形式列出，每条包含：负责人、具体任务、截止时间。如果会议记录中没有明确负责人或截止时间，请标注为「待确认」。

\verb|"""|

（粘贴会议记录）

\verb|"""|
\end{promptbox}

\subsection*{论文/报告核心观点提取}

\begin{promptbox}
请阅读以下报告，提取 3 到 5 个核心观点。每个观点用一句话总结，并标注出自报告的哪个部分。最后用 2 到 3 句话写一段整体摘要。

\verb|"""|

（粘贴报告内容）

\verb|"""|
\end{promptbox}

\section{信息提取与分类}

面对一堆杂乱的原始数据或文本，手动整理既枯燥又容易出错。AI 在这类「从非结构化文本中提取结构化信息」的任务上表现非常出色，几秒钟就能完成你手工要做半小时的分类和提取工作。

\subsection*{从产品评论中提取情感倾向}

\begin{promptbox}
请分析以下产品评论，将每条评论的情感分类为「正面」「负面」或「中性」，并用一句话说明判断理由。以表格形式输出，列名为：评论内容、情感分类、判断理由。

\verb|"""|

1. 这个产品质量太棒了，用了半年一点问题都没有！
2. 包装挺好看的，但实际用起来和描述不太一样。
3. 还行吧，没什么特别的。
4. 第二天就坏了，客服还不给退，差评！

\verb|"""|
\end{promptbox}

\subsection*{简历关键信息结构化提取}

\begin{promptbox}
请从以下简历中提取关键信息，以 JSON 格式输出，字段包括：姓名、联系方式、最高学历、毕业院校、工作年限、最近一份工作的职位和公司名称、核心技能（列表）。

\verb|"""|

（粘贴简历内容）

\verb|"""|
\end{promptbox}

\subsection*{客服工单自动分类}

\begin{promptbox}
以下是一批客服工单的问题摘要。请将它们自动分类到以下类别中：「产品质量」「物流配送」「退换货」「账户问题」「咨询建议」「其他」。以表格形式输出，列名为：工单编号、问题摘要、分类。

\verb|"""|

1. 收到的商品有破损
2. 订单显示已签收但我没收到
3. 想问一下你们的会员积分怎么用
4. 密码忘了，找回密码的邮件收不到
5. 买了两件，想退一件

\verb|"""|
\end{promptbox}

\section{文本转换与翻译}

翻译和改写是 AI 的传统强项。但「翻译」远不止是中英互译，它还包括在不同语言风格之间的转换：把学术语言变成大白话，把代码变成注释，把口语变成书面语。下面几个模板帮你搞定这些场景。

\subsection*{中英互译并保持语气}

\begin{promptbox}
请将以下中文翻译成英文。要求：保持原文的语气和风格（正式/口语化/幽默等），不要逐字直译，确保英文表达自然流畅。

\verb|"""|

（粘贴中文内容）

\verb|"""|
\end{promptbox}

\subsection*{把学术论文改写成科普文章}

\begin{promptbox}
以下是一篇学术论文的摘要。请将它改写成一篇面向普通读者的科普短文，要求：不使用专业术语（如果必须使用则加括号解释），语气轻松易懂，控制在 300 字以内。

\verb|"""|

（粘贴论文摘要）

\verb|"""|
\end{promptbox}

\subsection*{代码注释自动生成（代码辅助解析）}

\begin{promptbox}
请为以下代码添加详细的中文注释。要求：每个函数加上功能说明的文档注释，关键逻辑处加上行内注释，注释用中文书写，语言简洁明了。重点指出潜在的性能瓶颈或安全漏洞。

\verb|"""|

（粘贴代码）

\verb|"""|
\end{promptbox}

\begin{tipbox}
在进行代码或文本结构的重构、翻译时，如果你只需要复制最终的输出结果而不需要看 AI 啰里啰嗦的开场白（例如“好的，这是您的代码：”），建议在 Prompt 末尾加上一行「请直接输出结果本身，不需要任何多余的解释与寒暄」。这将极大提高复制粘贴的效率。
\end{tipbox}

\section{职场生存篇}

写周报、发邮件、做汇报……这些职场日常看似简单，却经常消耗你大量的时间和精力。用好以下模板，你可以把这些「文字活」的时间压缩到原来的几分之一，把精力留给更重要的事情。

\subsection*{一键生成周报}

\begin{promptbox}
你是一名经验丰富的职场人士。请根据以下工作要点，帮我写一份本周工作周报。格式要求：分为「本周完成」「进行中」「下周计划」三个板块，每个板块用编号列表，语气专业简洁。

本周做了这些事：
\begin{itemize}
  \item （列出你本周做的事情）
\end{itemize}
\end{promptbox}

\subsection*{润色年终总结}

\begin{promptbox}
请帮我润色以下年终工作总结。要求：保持原意不变，提升文字的专业度和可读性，语气沉稳大气，适当加入量化的表述来增强说服力。如果原文中有口语化的表达，请替换为更书面的说法。

\verb|"""|

（粘贴你的年终总结草稿）

\verb|"""|
\end{promptbox}

\subsection*{写一封得体的催款/拒绝邮件}

\textbf{催款邮件：}
\begin{promptbox}
请帮我写一封催款邮件。背景：对方是合作了两年的老客户，有一笔 5 万元的款项已经逾期 45 天。要求：语气礼貌但坚定，既要维护客户关系又要明确表达催款诉求，提及合同约定的付款周期，控制在 200 字以内。
\end{promptbox}

\textbf{婉拒邮件：}
\begin{promptbox}
请帮我写一封婉拒邮件。背景：一位同行邀请我参加下周的行业论坛做分享嘉宾，但我确实没有时间。要求：语气真诚表达感谢和歉意，给出合理的理由，留有未来合作的余地，控制在 150 字以内。
\end{promptbox}

\subsection*{长文档秒变大纲}

\begin{promptbox}
请阅读以下文档，将其内容整理成一份多层级的大纲。要求：一级标题不超过 7 个，每个一级标题下用要点列出核心内容，整体结构清晰、层次分明。

\verb|"""|

（粘贴文档内容）

\verb|"""|
\end{promptbox}

\section{内容创作篇}

如果你需要做自媒体运营、写营销文案或者制作短视频，AI 可以成为你最高效的「文案搭档」。它不一定能直接产出完美的成品，但绝对能帮你快速生成一个高质量的起点，让你在此基础上修改润色。

\subsection*{小红书种草文案}

\begin{promptbox}
你是一位小红书爆款文案写手。请为以下产品写一篇种草笔记：

产品：（产品名称和简介）

要求：
\begin{itemize}
  \item 标题要有吸引力，包含 emoji
  \item 正文分段清晰，每段不超过 3 行
  \item 要有真实感和个人体验感，不要太像广告
  \item 适当使用 emoji 增加阅读趣味
  \item 结尾加上 3 到 5 个相关话题标签
  \item 总字数 200 到 300 字
\end{itemize}
\end{promptbox}

\subsection*{短视频脚本规划}

\begin{promptbox}
请帮我写一个 60 秒短视频脚本。

主题：（你的视频主题）

要求：
\begin{itemize}
  \item 分为：开头hook（5秒抓住注意力）、正文内容（45秒）、结尾引导互动（10秒）
  \item 每个部分标注预计时长、画面描述、台词/旁白内容
  \item 语气口语化，适合真人出镜
  \item 开头第一句话要能让用户停下来不划走
\end{itemize}
\end{promptbox}

\subsection*{取标题的 5 种姿势}

\begin{promptbox}
请为以下文章内容提供 5 个不同风格的标题，并标注每个标题的风格类型：

1. 悬念型（引发好奇心）
2. 数字型（列出具体数量）
3. 痛点型（直击读者问题）
4. 对比型（制造反差感）
5. 利益型（突出读者收益）

文章内容概要：（简要描述你的文章内容）
\end{promptbox}

\section{生活与学习篇}

AI 不只是工作工具，在日常生活和学习中也能帮上大忙。从规划旅行到练习口语，从理解复杂概念到制定学习计划，以下模板覆盖了几个最实用的场景。

\subsection*{定制旅游攻略}

\begin{promptbox}
请帮我做一份旅游攻略。

目的地：（地点）

出行时间：（日期和天数）

预算：（总预算）

人数：（几个人、有没有老人小孩）

偏好：（美食/自然风光/历史文化/购物/其他）

要求：以日程表形式规划，每天包含上午、下午、晚上的安排，注明景点、交通方式、预计花费和餐饮推荐。
\end{promptbox}

\subsection*{AI 英语口语陪练}

\begin{promptbox}
你是一位耐心的英语口语老师。接下来我们进行一对一的英语对话练习。

规则：
\begin{itemize}
  \item 我用英语和你对话，你也用英语回复
  \item 如果我的表达有语法错误或不自然的地方，请在回复后用中文指出并给出更好的表达方式
  \item 用日常对话的语气，不要太学术
  \item 难度适中，适合中级英语学习者
\end{itemize}

我们从这个话题开始：今天的天气。请你先开始。
\end{promptbox}

\subsection*{复杂概念通俗化}

\begin{promptbox}
请用最通俗易懂的语言解释「量子力学」这个概念。要求：像讲给一个 10 岁小朋友听那样解释，可以用生活中的例子做类比，避免任何专业术语，控制在 200 字以内。
\end{promptbox}

\section{代码与技术篇}

即使你不是专业的程序员，AI 也能帮你写简单的脚本、排查代码错误、生成测试用例。下面的模板针对几个最常见的编程辅助场景，你只需要把自己的代码或需求填进去就能使用。

\subsection*{用自然语言生成代码}

\begin{promptbox}
请用 Python 写一个脚本，实现以下功能：

1. 读取当前目录下所有的 .csv 文件
2. 将每个文件中的「日期」列转换为统一的 YYYY-MM-DD 格式
3. 合并所有文件到一个总表
4. 按日期排序后输出为一个新的 merged.csv 文件

要求：添加必要的注释，使用 pandas 库，做好异常处理。
\end{promptbox}

\subsection*{AI 帮你 Debug}

\begin{promptbox}
以下 Python 代码运行时报错了，错误信息是：（粘贴错误信息）

请帮我：
1. 解释这个错误是什么意思
2. 找出代码中导致错误的具体位置
3. 给出修复后的完整代码
4. 解释修复的原因

\verb|"""|

（粘贴你的代码）

\verb|"""|
\end{promptbox}

\subsection*{自动生成单元测试}

\begin{promptbox}
请为以下 Python 函数编写单元测试。要求：使用 pytest 框架，覆盖正常输入、边界值和异常输入三种情况，每个测试用例加上清晰的注释说明测试目的。

\verb|"""|

（粘贴你的函数代码）

\verb|"""|
\end{promptbox}

\bigskip

以上模板只是一个起点。用的时候记得根据自己的实际需求调整细节，特别是背景信息、约束条件和输出格式。用得越多，你就越能感受到提示词的微小差异带来的巨大效果差别。

接下来的几章，我们会聊一些更实用的话题：多模态交互、翻车救援、安全与伦理，以及如何把零散的技巧整合成系统化的工作流。
